\chapter{Common Problems, Debugging, and Failure Analysis}
\label{chap:ch19}

\section{The Reality of Satellite Deep Learning Debugging}
In computer vision research, things rarely work on the first attempt. Training a multimodal, multi-scale transformer with adversarial loss involves dozens of moving parts. A single misaligned dimension, unnormalized channel, or inverted mask can corrupt months of computation.

This chapter presents a field-tested diagnostic guide to the 16 most common failure modes in multimodal satellite vision.

\section{The 16 Critical Failure Modes and Proven Remedies}

\subsection{1. Optical Radiometric Normalization Mismatch}
\begin{itemize}
    \item \textbf{Symptom:} Outputs are entirely pure white, washed out, or completely pitch black.
    \item \textbf{Root Cause:} Sentinel-2 images are stored as 16-bit integers ($[0, 10000]$). Dividing by 255 (standard 8-bit web image normalization) causes all reflectance values to exceed $1.0$ by a factor of 40!
    \item \textbf{Remedy:} Always normalize via $\text{clip}(\text{DN} / 10000.0, 0.0, 1.0)$.
\end{itemize}

\subsection{2. Swapped Channel Order}
\begin{itemize}
    \item \textbf{Symptom:} Model outputs surreal neon pinks or greens, and radar features are ignored.
    \item \textbf{Root Cause:} Python libraries handle channel orders differently (OpenCV uses BGR, PIL uses RGB, GeoTIFF uses arbitrary band indices). Slicing channels $0:2$ assuming they are SAR when they are actually Red/Green breaks the decoupled stem.
    \item \textbf{Remedy:} Enforce an explicit channel schema: \texttt{[VV, VH, B04, B03, B02, B08, Mask]}.
\end{itemize}

\subsection{3. Memory Layout / Dimension Transposition}
\begin{itemize}
    \item \textbf{Symptom:} \texttt{RuntimeError: Expected 4D tensor with shape [B, C, H, W] but got [B, H, W, C]}.
    \item \textbf{Root Cause:} NumPy and TensorFlow use channel-last ($[B, H, W, C]$); PyTorch requires channel-first ($[B, C, H, W]$).
    \item \textbf{Remedy:} Apply \texttt{tensor.permute(0, 3, 1, 2).contiguous()}.
\end{itemize}

\subsection{4. Inverted Cloud Mask}
\begin{itemize}
    \item \textbf{Symptom:} The final blended image contains clouds where clouds were, while clear areas are replaced by blurry neural predictions.
    \item \textbf{Root Cause:} Confusing $1 = \text{cloud}$ with $1 = \text{clear}$. In our pipeline:
    \begin{equation}
    \hat{\mathbf{Y}} = (1.0 - \mathbf{M}) \mathbf{X}_{S2} + \mathbf{M} \mathbf{P}_{\text{dir}}
    \end{equation}
    If $\mathbf{M} = 1$ on clear pixels, the formula discards the real clear pixels!
    \item \textbf{Remedy:} Always verify that $\mathbf{M} = 1$ corresponds to clouds and $0$ to clear sky.
\end{itemize}

\subsection{5. Misaligned SAR and Optical Imagery}
\begin{itemize}
    \item \textbf{Symptom:} Double coastlines, blurry roads, or ghosted city skylines.
    \item \textbf{Root Cause:} Sentinel-1 radar coordinates were not properly orthorectified with digital elevation models (DEM) or re-projected to match Sentinel-2's UTM grid.
    \item \textbf{Remedy:} Verify geocoding with sub-pixel cross-correlation before training.
\end{itemize}

\subsection{6. Noisy or Seasonally Shifted Reference Targets}
\begin{itemize}
    \item \textbf{Symptom:} Training loss stops decreasing; model reconstructs green grass where winter snow should be.
    \item \textbf{Root Cause:} The clear-sky target was acquired in summer while the cloudy image was taken in winter. The model is penalized for failing to predict summer crop phenology from winter radar data!
    \item \textbf{Remedy:} Use scale-decoupled $L_1$ pooling ($4\times 4, 8\times 8$) which reduces sensitivity to micro-temporal pixel phase variations.
\end{itemize}

\subsection{7. Over-Smoothed Outputs (Low Detail Ratio)}
\begin{itemize}
    \item \textbf{Symptom:} Cloud-cleared regions have acceptable PSNR but look like smooth, blurry watercolor paintings (Detail Ratio $< 0.50$).
    \item \textbf{Root Cause:} Solely relying on pixel-wise $L_1$/$L_2$ losses (the conditional mean trap).
    \item \textbf{Remedy:} Activate local standard deviation matching ($16\times 16$), multi-scale signed gradients, and the 11-channel high-pass PatchGAN.
\end{itemize}

\subsection{8. Global Color / Radiometric Drift}
\begin{itemize}
    \item \textbf{Symptom:} Entire reconstructed patches exhibit an unnatural yellow, purple, or dark cast.
    \item \textbf{Root Cause:} Final convolutional head bias initialized randomly with zero or Gaussian noise.
    \item \textbf{Remedy:} Initialize \texttt{final\_conv.bias} to \texttt{CLEAR\_CHANNEL\_MEANS} ($[0.0847, 0.0919, 0.1127, 0.1884]$).
\end{itemize}

\subsection{9. Unstable GAN Training and Mode Collapse}
\begin{itemize}
    \item \textbf{Symptom:} Discriminator loss drops to 0.0 within 50 steps; generator outputs catastrophic geometric checkerboard noise or repeating geometric artifacts.
    \item \textbf{Root Cause:} Discriminator is too strong, learning rates are unbalanced, or Lipschitz constraints are violated.
    \item \textbf{Remedy:} Apply Spectral Normalization to all discriminator layers, use mask-weighted Hinge loss, and reduce adversarial generator weight to $w_{\text{adv}} = 0.015$.
\end{itemize}

\subsection{10. Loss Plateaus Early in Training}
\begin{itemize}
    \item \textbf{Symptom:} Loss decreases slightly in epoch 1 and remains completely flat for 20 epochs.
    \item \textbf{Root Cause:} Learning rate too low or model trapped in a poor local minimum; radar features not propagating.
    \item \textbf{Remedy:} Implement a $2\times$ differential learning rate on the SAR stem ($2 \times 10^{-4}$ vs. $1 \times 10^{-4}$) and verify LayerScale initialization ($\gamma = 0.01$).
\end{itemize}

\subsection{11. Exploding or Vanishing Gradients}
\begin{itemize}
    \item \textbf{Symptom:} Loss outputs \texttt{NaN} (Not a Number) or weights freeze with zero updates.
    \item \textbf{Root Cause:} Extremely deep transformer attention layers without gradient clipping.
    \item \textbf{Remedy:} Apply \texttt{torch.nn.utils.clip\_grad\_norm\_(model.parameters(), max\_norm=1.0)} and Automatic Mixed Precision with \texttt{GradScaler}.
\end{itemize}

\subsection{12. Overfitting (Train Loss Falls, Val Loss Rises)}
\begin{itemize}
    \item \textbf{Symptom:} Train loss reaches $0.01$ while validation PSNR degrades from $30\text{ dB} \to 25\text{ dB}$.
    \item \textbf{Root Cause:} Model capacity is too high or data augmentation is absent.
    \item \textbf{Remedy:} Apply random horizontal/vertical flips, 90-degree rotations, synthetic opaque cloud core cutouts ($p=0.65$), and weight decay ($1\times 10^{-4}$).
\end{itemize}

\subsection{13. Spatial Data Leakage}
\begin{itemize}
    \item \textbf{Symptom:} Validation PSNR is artificially high ($38\text{ dB}$), but performance on a newly acquired external satellite tile is catastrophic ($19\text{ dB}$).
    \item \textbf{Root Cause:} Splitting dataset at the patch level rather than the geographical scene level.
    \item \textbf{Remedy:} Partition train, val, and test sets strictly by distinct ROI scene directories.
\end{itemize}

\subsection{14. Metric Calculation Range Errors}
\begin{itemize}
    \item \textbf{Symptom:} PSNR reported as $8.2\text{ dB}$ or negative values despite visually flawless reconstructions.
    \item \textbf{Root Cause:} Evaluating a $[0, 1]$ normalized image with a PSNR function that assumes $\text{MAX}_I = 255.0$ (or vice-versa).
    \item \textbf{Remedy:} In $10 \log_{10}(\text{MAX}^2 / \text{MSE})$, set $\text{MAX} = 1.0$ when inputs are in $[0, 1]$.
\end{itemize}

\subsection{15. Visible Perimeter Seams in Blended Outputs}
\begin{itemize}
    \item \textbf{Symptom:} Sharp, unnatural cookie-cutter borders surround reconstructed clouds.
    \item \textbf{Root Cause:} Blending with an unblurred binary mask ($\{0, 1\}$).
    \item \textbf{Remedy:} Apply an $11 \times 11$ Gaussian blur ($\sigma=2.0$) to the dilated cloud mask.
\end{itemize}

\subsection{16. Out of Memory (OOM) GPU Crashes}
\begin{itemize}
    \item \textbf{Symptom:} \texttt{CUDA out of memory. Tried to allocate 2.40 GiB}.
    \item \textbf{Root Cause:} High batch size or spatial self-attention memory explosion ($\mathcal{O}((HW)^2)$).
    \item \textbf{Remedy:} Reduce batch size to 8 or 16, enable AMP (\texttt{torch.amp.autocast}), and ensure transposed channel attention (MDTA) is used instead of spatial attention.
\end{itemize}

\section{Summary of Chapter 19}
\begin{enumerate}
    \item Normalization, channel ordering, and mask orientation must be verified with mathematical rigor before launching training.
    \item Over-smoothing is eliminated by moving away from naive $L_1$/$L_2$ losses to gradient, texture std, and high-pass PatchGAN losses.
    \item Mode collapse is prevented by Spectral Normalization and conservative adversarial weighting ($w_{\text{adv}} = 0.015$).
    \item Generalization requires scene-level dataset splitting to avoid spatial memorization.
\end{enumerate}

\section{Practice Exercises}
\begin{enumerate}
    \item If a training loss abruptly outputs \texttt{NaN} in epoch 3, what are the first two diagnostic actions you should take?
    \item How does calculating PSNR with $\text{MAX}=255.0$ on a $[0, 1]$ image distort the resulting decibel score?
    \item Why does a 3-step mask refinement (pool $\to$ dilate $\to$ Gaussian blur) resolve seam boundary artifacts?
\end{enumerate}

\subsection*{Solutions}
\begin{enumerate}
    \item First, verify whether gradient norm clipping (\texttt{clip\_grad\_norm\_}) was active, as unchecked gradient spikes produce infinities. Second, check whether any division by zero occurred (e.g., in cloud-masked losses without an $\epsilon = 10^{-8}$ denominator stabilizer) or if AMP \texttt{GradScaler} encountered severe overflow.
    \item If $\text{MAX}$ is incorrectly set to $255.0$ instead of $1.0$, the numerator inside the logarithm is $255^2 = 65{,}025$ times too large! The calculated PSNR will be artificially inflated by $+10 \log_{10}(65{,}025) \approx \mathbf{+48.13\text{ dB}}$!
    \item Dilation guarantees that cloud shadows and penumbra haze are encompassed within the mask, while Gaussian softening creates an alpha-blend transition over several pixels, eliminating step-function discontinuities.
\end{enumerate}
